% Actual cycle calculation supplied to maps/fourfolds for their D.2.44.
For the seven default semistable minimal rows with residual
$c=4$ or $8$, all primitive Mumford caps and ordinary clutch
$e=1$, the regular three-torus
$T=\langle e_1,e_2,\gamma\rangle$ generates
$H_3(X)=\mathbb Z/c$. The complex critical-curve cycle
has the following integral bounding-link values with $T$:
\[
 \begin{array}{c|c|c|c}
 \text{OS}&(a,b)&\text{zero index}&\text{critical linking}\\ \hline
 45&(1,8)&1,2,3,4&4,4,4,4\\
 55&(5,4)&2,3,4&37,7,27 .
 \end{array}
\]
In particular the latter three values are odd and the
former four are even. These are cycle computations;
no framing consequence is asserted here.

Here are the geometric markings used to make this an
actual cycle calculation. At each multiplicative site,
the recorded integral elliptic basis has determinant
one and takes the elliptic monodromy to its positive
transvection. In the resulting basis
$(u_1,u_2,\delta,\gamma)$, the short and long two-planes
are $(u_1,\delta)$ and $(u_2,\gamma)$, and the actual
valuation-period matrix is
\[
                    L=\begin{pmatrix}n&-bH\\aH&R\end{pmatrix}.
\]
These are the exact valuations, not reduced component
representatives.

At a rank-one site factor $L=\alpha r$, with $\alpha$
primitive. Choose a primitive short complement $\sigma$
with $\det(\alpha,\sigma)=1$ and a long basis
$(\beta,\ell)$ such that $r\beta=\rho>0$,
$r\ell=0$ and $\det(\beta,\ell)=1$.
Then $N\beta=\rho\alpha$. The vanishing elliptic
factor has its complex orientation $(\alpha,\beta)$,
so the transverse critical elliptic curve has
the complex orientation $(\sigma,\ell)$.
Its actual regular-fibre preimage is
$\sigma\wedge\ell$.
The marked primitive nodal-product model identifies
this torus with the critical elliptic curve by a
specialization chain contained entirely in that cap.
The determinant-one choices ensure that no covering
degree is introduced.

At the unique rank-two cusp, $\det L=1$.
Replace the long basis by
$(b_1,b_2)=(u_2,\gamma)L^{-1}$, so $Nb_j=a_j$ for
the short basis $(a_1,a_2)=(u_1,\delta)$.
Use the three integral toric support functions
\[
 \psi_1(r,s)=r(r-1)/2,\qquad
 \psi_2(r,s)=rs,\qquad
 \psi_3(r,s)=s(s-1)/2.
\]
Their degree matrix on the three complex-oriented
opposite-pair double curves, in the cyclic ray order
$(1,0),(0,1),(-1,1)$, is
\[
                     \begin{pmatrix}
                       0&-1&0\\1&1&-1\\-1&0&0
                     \end{pmatrix}.
\]
It has determinant $-1$.
Each entry follows from the actual wall formula
$\psi(v)+\psi(v')-\psi(w)-\psi(w')$.
The sum of the three curves therefore has degree
vector $(-1,1,-1)$.

On the regular torus the corresponding Chern classes
have mixed entries given by the affine slope
$\operatorname{Hess}(\psi)L$.
Consequently the regular-fibre class
\[
                    a_1\wedge b_1+b_1\wedge a_2
                                      +a_2\wedge b_2
\]
has exactly those three degree values.
It has zero image in the degree-two homology of
the dual torus, as does the sum of the double
curves: it has no long-long component.
The unimodular degree matrix and that fourth
coordinate therefore identify its specialization
with the literal sum of the three complex curves
integrally. This fixes the orientation and excludes
an undetected component in the cusp's fourth
homology generator.

Transport these five cap preimages to the common
regular fibre along the fixed base paths. The
transport chains lie outside the ordinary clutch
disc. Their sum is a fibre-vector representative
of the actual critical cycle in the full marked
presentation of D.2.41.
For example, at OS55, zero index $2$, the five
vectors in the ordered bivector basis
$(e_1e_2,e_1\delta,e_1\gamma,e_2\delta,e_2\gamma,\delta\gamma)$
are
\[
 \begin{gathered}
 (0,0,0,-12,0,5),\quad (3,0,1,-12,0,4),\\
 (-19,-23,-10,0,5,6),\quad
 (0,0,0,0,0,1),\quad (0,4,0,0,0,1).
 \end{gathered}
\]
Their sum is $(-16,-19,-9,-24,5,17)$.
The exact integral solution of the full relation
matrix $Mc=C_{\rm crit}$ has ordinary-clutch
loop vector $(0,0,37,0)$.

The linking extraction is the same actual Thom-chain
argument as in
\path{code/manifolds_regular_three_torus_proof.tex}.
Only the clutch-disc chains meet the chosen regular
fibre. A relation for $c$ times a Smith torsion
generator has clutch delta coefficient a unit
modulo $c$ in all seven rows, so $T$ generates
$H_3(X)$. A relation for the literal nullhomologous
critical cycle gives the displayed integer linking
values.

The complete lifts, Smith coefficient vectors and
clutch vectors are saved in
\path{records/manifolds_even_critical_cycles.json},
with executable
\path{code/manifolds_even_critical_cycles.py}.
The program also changes every short complement
and the cusp mixed preimage by an actual cap-kernel
class, and verifies that the resulting integer
linking is unchanged. It verifies that every
integral relation among the full presentation
columns has zero clutch delta coefficient.
Thus the integer does not depend on the chosen
Smith solution or those preimage choices.

The executable also repeats all seven calculations for ordinary
clutch $e=-1$, using its actual power vector $(0,0,0,1)$.
Each critical-linking integer is the negative of the corresponding
$e=1$ value, and every regular-three-torus linking coefficient
remains a unit modulo $c$. Thus the parity statements hold for
both simply connected ordinary-clutch signs.
